\section{Empirical quantiles, median, and interquartile range}

Quantiles were introduced for probability distributions through the cumulative
distribution function. For data, quantiles describe positions in the ordered
sample. They answer questions such as: where is the middle of the data? Where is
the lower quarter? Where is the upper quarter?

To discuss empirical quantiles, we first sort the data:
\[
x_{(1)}\leq x_{(2)}\leq \cdots \leq x_{(n)}.
\]
The parentheses indicate order statistics. The value \(x_{(1)}\) is the smallest
observation, and \(x_{(n)}\) is the largest observation.

\subsection*{Median}

The median is a value that divides the ordered data into two parts. Roughly half
of the observations lie at or below the median, and roughly half lie at or above
it.

For an odd sample size \(n=2m+1\), the median is usually taken as
\[
x_{(m+1)}.
\]
For an even sample size \(n=2m\), a common convention is to take the average of
the two middle observations:
\[
\frac{x_{(m)}+x_{(m+1)}}{2}.
\]

For example, for the ordered data
\[
1,\quad 2,\quad 2,\quad 3,\quad 100,
\]
the median is \(2\). For the ordered data
\[
1,\quad 2,\quad 2,\quad 3,
\]
a common median is
\[
\frac{2+2}{2}=2.
\]

\subsection*{Median versus mean}

The median and the mean both describe center, but they measure different kinds of
center. The mean is a balance point and is sensitive to extreme values. The
median is a positional center and is more resistant to extreme values.

Consider
\[
1,\quad 2,\quad 2,\quad 3,\quad 100.
\]
The mean is
\[
\bar{x}=\frac{108}{5}=21.6,
\]
while the median is \(2\). The value \(100\) pulls the mean strongly to the right,
but it does not move the median beyond the middle observation.

This does not mean that the median is always better. It means that the mean and
median answer different descriptive questions.

\subsection*{Empirical quantiles}

A \(p\)-quantile of a theoretical distribution is connected with the equation
\[
F_X(q_p)=p.
\]
For data, the empirical analogue is based on the empirical CDF. An empirical
\(p\)-quantile is a value below which approximately a proportion \(p\) of the data
lies.

There are several conventions for empirical quantiles, especially when
interpolation is used. Different software packages may give slightly different
answers for small datasets. The conceptual idea is stable: an empirical quantile
marks a probability position in the ordered sample.

\begin{remarkbox}
Empirical quantiles are not unique by one universal computational convention.
When exact numerical agreement matters, the convention used by the software or
textbook should be stated.
\end{remarkbox}

\subsection*{Quartiles and interquartile range}

The first quartile \(Q_1\) is an empirical \(0.25\)-quantile. The median \(Q_2\)
is an empirical \(0.5\)-quantile. The third quartile \(Q_3\) is an empirical
\(0.75\)-quantile.

The interquartile range is
\[
IQR=Q_3-Q_1.
\]
It measures the length of the central part of the data, ignoring the lowest and
highest quarters.

The IQR is often used as a robust measure of spread. It is less sensitive to
extreme observations than the standard deviation.

\subsection*{Relation to theoretical quantiles}

If a sample is generated from a distribution with theoretical \(p\)-quantile
\(q_p\), then the empirical \(p\)-quantile may be close to \(q_p\) for a large
sample. But the empirical quantile is computed from data, while the theoretical
quantile is defined by the probability law.

For example, in a standard normal model, the theoretical median is \(0\). A
finite sample from that model will usually have a sample median close to \(0\),
but not exactly equal to \(0\).

\begin{summarybox}
Empirical quantiles are positions in the ordered sample. They are data-based
counterparts of theoretical quantiles. The median describes positional center,
and the interquartile range describes the spread of the central part of the data.
\end{summarybox}
